\chapter{The Chart}
% OUTLINE Ch2 "The Chart (manifolds as the observer's coordinates)". Laws:
% CITATION (Riemann 1854/1868 posthumous via Dedekind; Whitney 1936; Lee as
% modern standard), LENS (no device appearance in this chapter -- dose spent
% at Ch1+Ch6), Vol-1 laws. The Riemann story at Cavendish grade. No metric
% anywhere -- the no-ruler theme continues. Gate: Kodo per-claim citation.

An observer lays down coordinates. That is not a mathematical event; it is a
bench event --- an origin chosen where the apparatus happens to sit, axes
aligned with the table's edges, units taken from the instrument at hand. A
second observer, in a second lab, lays down different ones. Neither set is
right, neither is wrong, and everything either of them publishes must
survive translation into the other's numbers, or it was an artifact of the
furniture and not a fact about the world. The physicist knows this
discipline in the bones: physics is what survives the change of observer.
This chapter is that discipline, written down as mathematics --- and, as in
the last chapter, the writing-down asks the reader to believe nothing they
do not already practice.

Begin with what coordinates actually are. A chart is a map --- in the
cartographer's sense before the mathematician's --- from a region of the
space under study onto a patch of ordinary number-space, one point to one
address, addresses varying continuously in the sense the last chapter
earned. The word ``map'' should be taken seriously: the chart is not the
territory, it is a labeled picture of a region, and like any honest map it
covers only its region. A space is called locally Euclidean when every point
has a neighborhood --- the last chapter's word, doing its work --- that some
chart covers. Nothing global is promised. The sphere of the Earth admits no
single flat map without tearing, and never needed one: an atlas of
overlapping regional maps has always sufficed the navigator, and it
suffices mathematics.

The load-bearing structure lives where two maps overlap. There, every point
carries two addresses, and the rule converting one address into the other
--- the transition map --- is a map from number-patch to number-patch:
observer-to-observer translation, expressed entirely in coordinates. Now
the construction takes its one step, and the reader should watch how little
it costs. Demand that the transitions be continuous, and the patchwork is a
topological manifold. Demand more --- that the transitions be
differentiable, as maps of number-space, where differentiation already has
its classical meaning --- and the patchwork is a smooth manifold. Notice
what the demand touches: not the space, the \emph{translations}. Smoothness
is imposed on the observer-to-observer dictionaries, which live in ordinary
number-space where calculus has worked since Newton; the space itself then
inherits a calculus through its charts, region by region, with the
transition maps guaranteeing that what one observer computes another can
re-derive. Differential geometry begins as a treaty among observers.

The idea has a birthday, and the story deserves telling at the grade this
series tells its instrument-stories. In 1854 in G\"ottingen, Bernhard
Riemann, twenty-seven and a candidate for the right to lecture, submitted
--- as custom required --- three topics for his trial lecture, expecting,
as custom also all but promised, that the faculty would take the first. His
examiner was Gauss, and Gauss, against custom, chose the third: the
foundations of geometry. What Riemann produced under that surprise, in a
few pressed weeks, was the lecture ``On the Hypotheses Which Lie at the
Foundations of Geometry'' --- delivered the tenth of June, 1854, to a
faculty audience mostly unequipped to follow it, with famously almost no
formulas, and by the account that has come down to us it moved Gauss, who
was not a man reports describe as easily moved. Riemann's proposal was the
one this chapter has been building: space as an $n$-fold extended
magnitude, charted region by region, its deeper structure --- the metric
relations, the lengths --- not given by logic in advance but to be
determined, as he said, by experience. The lecture was published only in
1868, two years after his death, by Dedekind; the physicists took the rest
of a century to spend what it contained. The reader should mark the
physical lens Riemann himself ground: the question of which geometry space
wears, he insisted, belongs to measurement, not to axiom. This series
agrees, and holds the thought for later volumes.

Between Riemann's lecture and the modern textbook stands one more name. The
patchwork definition --- charts, atlases, smooth transitions --- was made
fully precise by Hassler Whitney in 1936, and Whitney proved in the same
stroke the theorem that keeps the abstraction honest: every smooth
$n$-dimensional manifold, however abstractly assembled from patches, can be
embedded in an ordinary Euclidean space of dimension $2n+1$ --- realized as
an honest surface sitting in a big enough number-space. The abstract view
costs nothing, then: whatever a patchwork manifold is, it is never anything
more exotic than a smooth surface seen without reference to the room it
sits in. The assembled modern treatment is standard in the literature ---
Lee's \emph{Introduction to Smooth Manifolds} is the reference this series
will cite for the details it declines to repeat --- and nothing in it goes
beyond the two demands already stated: charts that cover, transitions that
are smooth.

So what, in the end, \emph{is} the manifold, if every chart is somebody's
furniture? It is the residue --- the content that survives every admissible
change of chart, the facts on which all observers' translations agree. That
is not a poetic gloss; it is the working definition the construction just
earned, and it is the mathematical spelling of the discipline the chapter
opened with: physics is what survives the change of observer, and a
manifold is what survives the change of chart. The reader should also
notice, one more time, what is still absent. A manifold, even a smooth one,
has no lengths. No distances, no angles, no areas; nothing so far prefers
one chart's scale to another's. The ruler is still chapters away, and the
amount of structure standing without it continues to grow.

One thing a chart does supply, at every point it covers, is the possibility
of asking a pointed question: standing at this point, at this instant, in
which direction --- and how fast? The physicist's word for the answer is
velocity, and velocity does not need the whole space to be defined; it
needs only an instant's worth of room around one point. What velocities are,
what lives beside them, and how the two great families of components the
physicist writes --- the raised indices and the lowered ones --- arise from
one construction, is the next chapter's subject, and it is where the second
of this book's two title instruments, linear algebra, enters the shop.
